\section*{第 \ref{ch:g6:producing-our-food} 章——生产我们的食物}

\begin{solution}{exo:g6:producing-our-food:1}
\omterm{def:g6:producing-our-food:farming}{作物种植}管理的是\omterm{def:g5:food-webs:roles}{生产者}——即为获取其\omterm{def:g6:plants-are-producers:matter}{有机物}而种植的\omterm{def:g1:plants-around-us:plant}{植物}。
\omterm{def:g6:producing-our-food:farming}{畜牧养殖}管理的是\omterm{def:g5:food-webs:roles}{消费者}——即为获取其产品而饲养的\omterm{def:g1:animals-around-us:animal}{动物}：
\omterm{def:g2:animals-grow-up:born}{乳汁}、肉、蛋、羊毛。
\end{solution}

\begin{solution}{exo:g6:producing-our-food:2}
谷粒：小麦植株的\omterm{def:g2:from-seed-to-plant:seed}{种子}（即其打包好的储备）。
马铃薯：马铃薯植株的\omterm{ex:g4:plants-without-seeds:bulbtuber}{块茎}。苹果：苹果树的\omterm{prop:g3:flowers-fruits-seeds:transform}{果实}。
牛奶：奶牛的哺育产物。卵：母鸡的
打包好的育雏室。
\end{solution}

\begin{solution}{exo:g6:producing-our-food:3}
通过\omterm{def:g6:producing-our-food:fermentation}{微生物}取食实现的原料转化。
面包（原料：面粉面团，工作者：\omterm{ex:g6:producing-our-food:bread}{酵母}）与
酸奶（原料：牛奶，工作者：使牛奶变酸的\omterm{def:g6:producing-our-food:fermentation}{微生物}）；
奶酪和醋也属此类。
\end{solution}

\begin{solution}{exo:g6:producing-our-food:4}
\omterm{ex:g6:producing-our-food:bread}{酵母}——一种单\omterm{def:g5:first-look-at-cells:cell}{细胞}\omterm{def:g6:decomposers-and-soil:decomposers}{真菌}——以面粉中的糖类为食，
并释放出气体；有弹性的面团把气体截留成气泡，使面团膨胀。
烤箱的热量终止了\omterm{ex:g6:producing-our-food:bread}{酵母}的工作，
使膨胀的面包瓤定型，连同其中的孔洞一起。
\end{solution}

\begin{solution}{exo:g6:producing-our-food:5}
\omterm{def:g6:producing-our-food:fermentation}{微生物}取食牛奶中的乳糖，释放出酸，
酸使牛奶增稠凝固，
而清新的酸味正是这种酸本身。
\end{solution}

\begin{solution}{exo:g6:producing-our-food:6}
腌制、冷藏（冰箱）、干燥，还有烹饪：
每种方法都使有害\omterm{def:g6:producing-our-food:fermentation}{微生物}失去所需的\omterm{def:g6:exploring-our-environment:conditions}{条件}（水分、\omterm{def:g6:exploring-our-environment:conditions}{温度}），
正如低温会减缓堆肥的腐熟。
\end{solution}

\begin{solution}{exo:g6:producing-our-food:7}
水：农民依赖（有时也引入）降雨。
光：行距株距适当，使每株都分得足够的光。
温暖：选在暖季播种。空气：
开阔的地面，松散的\omterm{def:g6:decomposers-and-soil:soil}{土壤}。\omterm{prop:g4:what-plants-need:needs}{无机盐}：
收获后施用粪肥，补充养分。
\end{solution}

\begin{solution}{exo:g6:producing-our-food:8}
(1) \omterm{def:g6:exploring-our-environment:components}{生物}体与阶段：牛奶——\omterm{def:g5:food-webs:roles}{消费者}的产品。
(2) 循环中的位置：\omterm{def:g5:food-webs:roles}{消费者}的产品，经\omterm{def:g6:producing-our-food:fermentation}{微生物}加工。
(3) 原料为牛奶；工作者使其变酸、凝结成凝乳，再将压制的奶酪熟化数月——
孔洞、外皮和气味都是它们的记录。
(4) 追溯路径：
牛奶来自奶牛，奶牛吃草，草依赖光。
\end{solution}

\begin{solution}{exo:g6:producing-our-food:9}
\omterm{def:g6:producing-our-food:fermentation}{微生物}是有\omterm{def:g6:exploring-our-environment:conditions}{条件}限制的活工作者：寒冷几乎使其停止取食。
冰箱的低温保护了食物（使腐败菌停止活动）；
冬季堆肥中，低温延缓了堆肥（使\omterm{def:g6:decomposers-and-soil:decomposers}{分解者}停止活动）。
原理：
分解与变质的快慢，
取决于\omterm{def:g6:exploring-our-environment:conditions}{条件}所允许的程度。
\end{solution}

\begin{solution}{exo:g6:producing-our-food:10}
“\omterm{def:g6:producing-our-food:fermentation}{微生物}”一词指的是一大批形形色色的微小\omterm{def:g6:exploring-our-environment:components}{生物}。
其中一些出现在不当之处，会引发疾病——这些正是洗手要防范的；
另一些在受控\omterm{def:g6:exploring-our-environment:conditions}{条件}下被选育和管理，能使面包\omterm{def:g6:producing-our-food:fermentation}{发酵}、牛奶凝结。
因此，是敌是友，
并不取决于\omterm{def:g6:producing-our-food:fermentation}{微生物}的本性，
而取决于\omterm{def:g5:biodiversity:species}{物种}、场所与管理方式是否匹配。
\end{solution}

\begin{solution}{exo:g6:producing-our-food:11}
火腿是经由猪转化而来的，其体重所耗的饲料是猪自身重量的许多倍，
另加猪自身消耗的那部分，以及占用的土地和水——
这是\omterm{def:g5:food-webs:roles}{消费者}层级的代价。面包中的小麦只经过一个层级就来自光。
浪费一个三明治，两笔账都会被浪费；
火腿的那笔账更长。
\end{solution}

\begin{solution}{exo:g6:producing-our-food:12}
他们管理的是工作者的\emph{\omterm{def:g6:exploring-our-environment:conditions}{条件}}：面团角落的\omterm{def:g6:exploring-our-environment:conditions}{温度}、
保存并定期喂养的酵头，
以及盐量与时机，这些都经过世代反复试验调整。
配方之所以有效，
是因为保持正确的\omterm{def:g6:exploring-our-environment:conditions}{条件}，就能让正确的\omterm{def:g6:exploring-our-environment:components}{生物}繁盛，
无论人们是否知道这些\omterm{def:g6:exploring-our-environment:components}{生物}的存在。他们在看不见\omterm{def:g6:producing-our-food:fermentation}{微生物}的情况下培育\omterm{def:g6:producing-our-food:fermentation}{微生物}，
而且做得很好。
\end{solution}

\begin{solution}{pb:g6:producing-our-food:1}
\textbf{1.} 每粒谷物都是一颗\omterm{def:g2:from-seed-to-plant:seed}{种子}：里面包着一株微小的小麦\omterm{prop:g2:from-seed-to-plant:cycle}{幼苗}和
一份养料储备，是母株为\omterm{prop:g2:from-seed-to-plant:cycle}{幼苗}最初几天准备的——
并非为我们而准备，尽管我们把它挪作了他用。

\textbf{2.} 在小麦植株自身的绿叶中进行，
以水、\omterm{def:g6:decomposers-and-soil:soil}{土壤}中的\omterm{prop:g4:what-plants-need:needs}{无机盐}和空气中的二氧化碳为原料，
以光为\omterm{prop:g3:balanced-diet:jobs}{能量}来源。

\textbf{3.} 播种后，谷粒经历其\omterm{def:g2:born-grow-die:cycle}{生命周期}：
\omterm{def:g2:from-seed-to-plant:germination}{萌发}、长成植株，再结出新谷粒。出售并磨成面粉后，
其\omterm{def:g2:born-grow-die:cycle}{生命周期}即告终结，养料储备便归了我们——
这份本为\omterm{prop:g2:from-seed-to-plant:cycle}{幼苗}准备的午餐，被人吃掉了。

\textbf{4.} 借用什么，就归还什么：粪肥腐熟后
补充了收获时被带走的\omterm{prop:g4:what-plants-need:needs}{无机盐}——
满足了对\omterm{prop:g4:what-plants-need:needs}{无机盐}的需求。

\textbf{5.} \omterm{ex:g6:producing-our-food:bread}{酵母}是活的工作者：温暖时，它
取食快、干活快；天冷时，它便闲置——堆肥工作队的规则，
在面包房的规模上同样适用。

\textbf{6.} 它以面粉中的糖类为食，并释放出气体；
有弹性的面团把气体截留成气泡，气泡
使面团体积翻了一番。

\textbf{7.} 寒冷的时段使\omterm{ex:g6:producing-our-food:bread}{酵母}闲置——几乎没有取食，
几乎不产生气体。两组“工作队”无意中完成了一次公平的对照实验，
探究\omterm{def:g6:exploring-our-environment:conditions}{温度}的影响：面团相同、当天相同，只有一个\omterm{def:g6:exploring-our-environment:conditions}{条件}不同，
两次\omterm{def:g6:producing-our-food:fermentation}{发酵}可供比较。

\textbf{8.} 高温杀死了\omterm{ex:g6:producing-our-food:bread}{酵母}，结束了
这一班次；定型后的面包屑保留着气泡留下的孔洞——
这支工作队留下的印记，在烘烤中定格。

\textbf{9.} 切片来自面包，面包来自面团，面团来自
面粉，面粉来自谷物，谷物来自小麦植株的
叶片——在那里，以无机原料借助光制造\omterm{def:g6:plants-are-producers:matter}{有机物}：
这条链条的源头，是七月田野里的阳光。

\textbf{10.} 切片：一层（从\omterm{def:g5:food-webs:roles}{生产者}直达餐桌，再加一道
\omterm{def:g6:producing-our-food:fermentation}{微生物}加工）。黄油：两层（草到牛，再到
奶油）。火腿：两层，代价最大——猪
在变成火腿之前，已燃烧并排出了饲料中的大部分物质。
所以切片在物质上最省；火腿
则代价最高。

\textbf{11.} 堆肥中的\omterm{def:g6:decomposers-and-soil:decomposers}{分解者}——蚯蚓、潮虫、
\omterm{def:g6:decomposers-and-soil:decomposers}{真菌}和看不见的工作队——把面团中的\omterm{def:g6:plants-are-producers:matter}{有机物}
分解为\omterm{def:g6:decomposers-and-soil:soil}{腐殖质}和\omterm{prop:g4:what-plants-need:needs}{无机盐}：从食物到来年\omterm{def:g6:decomposers-and-soil:soil}{土壤}肥力，
只差一轮循环。

\textbf{12.} 例如：“我们的面包由一个
\omterm{def:g5:food-webs:roles}{生产者}利用\omterm{def:g6:plants-are-producers:matter}{无机物}质和光制成，由\omterm{def:g6:producing-our-food:fermentation}{微生物}
工作队\omterm{def:g6:producing-our-food:fermentation}{发酵}，完全不需要\omterm{def:g5:food-webs:roles}{消费者}这一层——
不过隔壁的火腿三明治却需要一层。”
\end{solution}
